\documentclass[11pt,a4paper]{article}
\usepackage[T1]{fontenc}
\usepackage[utf8]{inputenc}
\usepackage[english]{babel}
\usepackage{lmodern}
\usepackage[margin=27mm]{geometry}
\usepackage{microtype}
\usepackage[hidelinks]{hyperref}
\title{Energy, Latency, and Memory Cost of Continuous RGB and LWIR Perception on Embedded Devices}
\author{Mateusz Pomianek}
\date{Version 0.2.0 -- 21 November 2026}
\begin{document}
\maketitle

\begin{abstract}
Embedded multispectral perception must be evaluated using instrumented end-to-end cost rather than nominal accelerator throughput. What is the measured energy cost per complete annotated frame for fixed RGB, LWIR and fused perception under predeclared quality constraints? The prospective design is as follows: Measure sensor capture, preprocessing, inference, fusion and reporting on a documented target device with actual power instrumentation, thermal stabilization and repeatable workloads. The primary outcome is Measured joules per complete processed frame among configurations meeting fixed recall, false-alarm and latency bounds. The manuscript defines a research question and an evaluation strategy; it does not report experimental estimates.
\end{abstract}

\section{Introduction}
Embedded multispectral perception must be evaluated using instrumented end-to-end cost rather than nominal accelerator throughput. The scientific contribution under consideration is a falsifiable comparison in which alternative explanations, measurement limitations and distinct hypotheses from related ARC Robo studies can be assessed independently. Both the source of observations and the assumptions needed to interpret results must be documented before any confirmatory comparison.

\section{Detailed scientific scope and novelty boundary}
The research compares actual embedded system resource costs of continuous RGB, LWIR and fixed dual-sensor detection. The complete pipeline includes sensing, preprocessing, timing/registration, inference, fusion and reporting. Nominal TOPS, desktop benchmark time and GPU utilization are not physical measurements of device energy. P03 provides a fixed-mode reference for optional adaptive switching in P07.

\section{Falsifiable hypotheses}
\paragraph{H1.} Measured energy per complete processed frame differs substantially from estimates obtained from inference-only latency or nominal accelerator specifications.

\paragraph{H2.} Adding sensor capture, alignment and thermal constraints changes the feasible quality-energy-latency Pareto frontier relative to an inference-only comparison.

\paragraph{H3.} Sustained continuous operation increases upper-tail latency or energy per accepted frame relative to short warm-up runs under at least one predeclared fixed pipeline condition.

\section{Operationalization and study population}
\subsection{Independent unit of analysis}
An independent device-and-configuration workload run; consecutive frames from one run are clustered.

\subsection{Data, sensors and ground-truth requirements}
Named embedded device, camera and compute configuration, independent physical power meter, idle/reference traces, acquisition and processing timestamps, memory logs, detector version and legally usable labelled scene workloads. Data and models will be versioned with checksums, licensing information, instrument parameters, selection rules and clock provenance. The held-out test cohort must be independent at the acquisition-session, scene or participant level rather than at the individual frame level.

\subsection{Primary and secondary endpoints}
The primary estimand is defined as follows: Instrumented Joules per complete accepted frame for a configuration satisfying predeclared minimum recall, false-positive-per-frame and latency constraints.

The secondary measurements will be: Power distribution, input-to-output p50/p95/p99 delay, peak RAM/flash, sustained FPS, thermal state, drop fraction and number of infeasible conditions. The exact numerator, denominator, outcome threshold and permitted exclusions must be fixed in the protocol before inferential analysis.

\section{Planned experiments and research sequence}
\subsection{A. Instrumentation calibration and accounting boundary}
Select a specified embedded target, compatible sensor modes and a physical power instrument with documented sample rate and placement. Measure idle baseline, warm-up, thermal and supply-state effects; lock pipeline boundaries and input resolution.

\subsection{B. Repeatable fixed-mode measurement}
Run matched labelled scene workloads with fixed RGB, fixed LWIR and fixed dual-stream processing. Randomize configuration order and repeat sessions, logging acquisition-to-output timestamps, peak RAM/flash, power traces and disabled configurations.

\subsection{C. Sustained-load feasibility test}
Repeat long enough to identify thermal throttling, memory pressure, bandwidth effects and processing deadline misses. Report p50/p95/p99 latency, throughput and energy integrated over the full workload, not only model inference.

\subsection{D. Quality-constrained operating envelope}
Independently validate recall and false positives at frozen thresholds and identify configurations satisfying minimum detection quality and latency. Compute energy cost only for feasible configurations while explicitly listing failures to meet constraints.

\section{Comparative design and statistical inference}
Integrate physical power with explicit idle correction and instrumentation uncertainty. Compare paired independent runs; report device-run confidence intervals and detection-quality estimates on separate held-out scenes. Model warm-up and thermal changes rather than discarding them. When human participants or independent trajectories are needed, their number must follow a documented power simulation or precision calculation after a feasible pilot. Report uncertainty, negative results and sensitivity to failure strata. Thresholds, model configurations and exclusion criteria must be frozen before final testing.

\section{Evidence gates, reproducibility and ethical boundary}
Physical target-device power measurement is required. No energy claims from host-side processing timings, simulator performance or marketing TOPS values. Each study phase will preserve source SHA or release, configuration, labels, independent oracle, random seeds, raw event timing, and an auditable trace from inputs to computed figures. Simulated evidence is distinct from measured hardware outcomes. Participant studies require prior ethical review and consent.

\section{Results}
Empirical results, tables and uncertainty intervals will be incorporated after execution of the documented evaluation protocol and independent verification of the input evidence.

\section{Discussion}
The planned comparison is intended to determine whether the proposed mechanism yields an identifiable gain over the stated controls. Interpretation must distinguish failure of a scientific hypothesis from an incomplete implementation, and should describe practical operating conditions under which conclusions remain valid.

\section{Conclusion}
This study defines a testable research problem, matched comparators and a primary evaluation target. A substantive empirical conclusion remains conditional on reproducible measurements and appropriate methodological review.

\section{Data and code availability}
The authoritative scientific source is maintained in the \texttt{MatPomGit/arc-rgbt} repository. Dataset permissions, input manifests, source commits, evaluation scripts and permitted public artifacts will be identified in the final version. No participant records or proprietary instrumentation details are included here.
\end{document}
